\documentclass{article}

\usepackage{caption}
\usepackage{datatool}
\usepackage{etoolbox}
\usepackage{expl3} 
\usepackage[colorlinks=true, urlcolor=blue, linkcolor=blue]{hyperref}
\usepackage{xstring}
\usepackage{verbatim}

\newcommand\Thema{Armee}
\newcommand\city{Ulm}
\newcommand\datum{11.05.2026}
\newcommand\timeVorb{07:30-09:00}
\newcommand\timeEinf{09:00-09:30}
\newcommand\timeFrakOne{09:30-11:00}
\newcommand\timePauseOne{11:00-11:15}
\newcommand\timeAuss{11:15-12:30}
\newcommand\timeMittag{12:30-13:00}
\newcommand\timeFrakTwo{13:00-13:30}
\newcommand\timeFinVerh{13:30-14:00}
\newcommand\timePlenar{14:00-15:15}
\newcommand\timeDebr{15:15-15:30}
\newcommand\politiker{Andrea Wechsler}
\newcommand\politikerOffice{Mitglied des Europäischen Parlaments}
\newcommand\stadtvertreter{Nalan Schmidt}
\newcommand\stadtvertreterOffice{Anlauf- und Koordinierungsstelle Bürgerdialog (Open Government) der Stadt Ulm}
\newcommand\localSupport{das Europe Direct Ulm}
\newcommand\sponsor{die Vertretung der Europäischen Kommission in München}
\newcommand\jefvorsitz{Farras Fathi}
\newcommand\gendervorsitz{Landesvorsitzender}
\newcommand\evpLeader{Tobias}
\newcommand\evpRoom{Kleiner Sitzungssaal}
\newcommand\sdLeader{Antonia}
\newcommand\sdRoom{Veranstaltungsraum Europe Direct}
\newcommand\reLeader{Alica}
\newcommand\reRoom{Besprechungszimmer Donaubüro I}
\newcommand\fifthLeader{Farras}
\newcommand\fifthRoom{Besprechungszimmer Öffentlichkeitsarbeit}
\newcommand\pfeLeader{Hannah}
\newcommand\pfeRoom{Besprechungszimmer BM III}
\newcommand\numSuS{27}
\newcommand\location{im Ulmer Rathaus}
\newcommand\fifthGroup{Linke}
\newcommand\anzahlcomm{3}
\newcommand\totcoms{26}
\newcommand\minmems{25}
\newcommand\maxmems{90}




\title{SimEP How-To Guide}
\author{Svea \& Vincent}

\newcommand{\KOM}{Vincent}
\newcommand{\EP}{Svea}
\date{April 2026}


\IfEqCase{\Thema}{%
	{Green Deal}{\newcommand{\THEMA}{zur Europäischen Klimapolitik im Rahmen des Green Deal}}%
	{Asyl}{\newcommand{\THEMA}{zu einem gemeinsamen Europäischen Asylsystem}}%
	{Armee}{\newcommand{\THEMA}{zu einer Europäischen Armee}}
}

\ExplSyntaxOn
\cs_new:Npn \jef_strip_prefix:n #1
	{
		\tl_set:Ne \l_tmpa_tl {#1}
		\regex_replace_once:nnN
			{\A\s*(?:(?:im|am)\s+|(?:in|an)\s+(?:der|dem|den)\s+|(?:der|die|das|den|dem|ein|eine|einen|einem|einer)\s+)}
			{}
			\l_tmpa_tl
		\tl_use:N \l_tmpa_tl
	}
\NewDocumentCommand{\stripPrefix}{m}{\jef_strip_prefix:n{#1}}
\ExplSyntaxOff
 
 

\begin{document}

	\section{Allgemeines}
	
	\begin{itemize}
	\item Bitte wenn möglich mit dem Deutschlandticket fahren, falls vorhanden, es gelten die Erstattungsregeln der JEF Bayern
	\item Weitere Informationen zu eurer Vorbereitung findet ihr \href{https://drive.google.com/drive/folders/1zBLLjXlWXjD55VzzwxbWSNNgrS0o0mv4}{hier}
	\end{itemize}
	
	\subsection{Grober Teamer-Zeitplan}	
	
	\begin{itemize}
	\item Vorabend der SimEP: gemeinsames Abendessen und letzte Absprachen (ist natürlich nicht verpflichtend, aber halt kostenlos)
	\item Tag der SimEP: 07.30 Uhr-16.00 Uhr SimEP; danach kurze Nachbesprechung
	\item Abreise am besten ab 17.30/18.00 Uhr
	\end{itemize}
	
	\subsection{Teamer:innen-Rollen}
	
	\begin{itemize}
	\item Foto-Beauftrage:r
	\item Social-Media-Beauftragte:r + Berichterstattung
	\item Fraktionsleitung + Assistenz
	\item Ausschussleitung + Assistenz
	\end{itemize}
	
	Die Lehrkräfte sollten sich normalerweise nicht in den Fraktions-/Ausschussräumen aufhalten, um die Schülerinnen und Schüler nicht abzulenken. Es kann natürlich Ausnahmen geben, Alltägliches aus der Schule sollte aber nicht Teil des Planspiels sein!

	
	\subsection{Zugangsdaten}
	
	\begin{center}
	\begin{tabular}{c|c|c}
		Website & Benutzername & Passwort \\
        \hline
        Antragsgrün	& semerak@jef-bayern.de	& 1DD1smln \\
        Slido 		& simep@jef-bayern.de	& Metsola24! \\
	\end{tabular}
	\end{center}	
	
	\subsection{Phasen-Übersicht}
	
	\begin{tabular}{cl|cl}
		\timeVorb: 	 	& Vorbereitung 			& \timeFrakTwo:	& 2. Fraktionssitzung \\
		\timeEinf: 	 	& Einführung 			& \timeFinVerh:	& Finale Verhandlungsphase \\
		\timeFrakOne:	& 1. Fraktionssitzung 	& \timePlenar:	& Plenarsitzung \\
		\timeAuss: 		& Ausschusssitzung 		& \timeDebr:	& Debriefing \& Feedback \\
		\timeMittag: 	& Mittagspause 			& 
	\end{tabular}		
	
	\newpage	
	
	
	\section{Ablauf}
	\subsection{Vorbereitung (\timeVorb)}
	\begin{itemize}
		\item Raumtour \& Raumschilder aushängen
        \item Technikcheck in Fraktionsräumen [Fraktionsteamer:innen]
		\item Technikcheck Plenum [\KOM\ \& \EP]
		\begin{itemize}
			\item fester Laptop
			\item Standby-Modus aus
			\item alle (und nur) nötigen Programme offen \& laufen
		\end{itemize}
		\item Dekoration Plenum \& Mappen abzählen [Fraktionsteamer:innen]
		\item Schüler in Empfang nehmen [\KOM\ \& \EP]
%		\item \politiker\ in Empfang nehmen [\KOM\ \& \EP]		
%		\item \stadtvertreter\ in Empfang nehmen [\KOM\ \& \EP]
	\end{itemize}
	
	\subsection{Einführung (\timeEinf)}
	\begin{itemize}
%		\item Kurze Begrüßung/Intro [\KOM]
%		\item Grußwort [\politiker]
%		\item Grußwort [\stadtvertreter]
		\item Präsentation Briefing [Hannah \& \KOM]
        \item SuS-Anzahl nachzählen \& Mappen entsprechend anpassen [Fraktionsteamer:innen]
        \item Aufgaben:
        \begin{itemize}
            \item Fraktionsleiter am Podium
            \item Rest verteilt Mappen auf Stichwort \newline
            $\to$ Mit kleinster Fraktion anfangen \newline
            $\to$ Auf Geschlechterverteilung achten
        \end{itemize}
    \end{itemize}
	
	\subsection{1. Fraktionssitzung (\timeFrakOne)}
	\begin{enumerate}
		\item Icebreaker (Bildkarten) [5 min]: \newline Assoziationen mit Fraktion \& Thema \newline
		$\to$ am besten in lockerer Atmosphäre (bspw. im Kreis stehen statt sitzen)
		\item Vorstellung \& Besprechung der Unterlagen, insbesondere den Gesetzentwurf durchgehen \newline $\to$
		Jeden Abschnitt des Entwurfs grob durchgehen, sodass die SuS einen Überblick über den gesamten Entwurf haben [15 min]
		\item Restliche Präsentation [10 min]
		\begin{itemize}
			\item Rolle der Schüler deutlich machen: \newline
            Nicht persönliche Meinung vertreten, aber auch kein Fraktionszwang
			\item Überblick über Fraktion
			\item Weltanschauung der Fraktion
			\item Rolle der Nationalität/Länderpapiere
			\item Weiteren Ablauf des Gesetzgebungsprozesses/der SimEP klarstellen
            \item Raum für Fragen
		\end{itemize}
		\item Slido-Quiz über den Entwurf und die Position der Fraktion \THEMA [10 min]
        \item Zufällig in \anzahlcomm\ gleichgroße Kleingruppen einteilen \newline $\to$
        Jeder Kleingruppe durch Namensschilder Ausschuss zuweisen (einen der Namensschilderbögen geben)
		\item In den Auschussgruppen [50 min]
		\begin{enumerate}
			\item Gesetzesentwurf (besonders zugewiesener Abschnitt) lesen
			\item EINEN (!) ÄA pro Ausschussgruppe erarbeiten \newline (muss sich auf einen inhaltlichen Aspekt beschränken)
	        \item Ausschusssprecher:in festlegen		
            \item ÄA in Antragsgrün einbringen (Beispiele \href{https://drive.google.com/drive/folders/14sA-kBBqEivc-man3TKYYKYBKKZzBpAR}{hier})
            \begin{itemize}
            		\item Entsprechenden Absatz auswählen
            		\item Änderungen vornehmen (\textit{= zu Ersetzendes/Entfernendes löschen bzw. Neues im Stil des Textes einfügen})
            		\item Nach ganz unten scrollen
            		\begin{itemize}
					\item Unter \textbf{Name} die \textit{Fraktion} eintragen
            			\item Unter \textbf{Gremium,LAG...} den betreffenden \textit{Ausschuss} eintragen
            			\item Unter \textbf{E-Mail} \textit{simep@jef-bayern.de} eintragen
            		\end{itemize}
            		\item Bei ÄA bestätigen auf \textit{Veröffentlichen} klicken
            \end{itemize}
		\end{enumerate} 
	\end{enumerate}
	
	\subsection{Zwischenpause (\timePauseOne)}
	\begin{itemize}
		\item Reihenfolge ÄA festlegen (weitreichendste zuerst) [Ausschussteamer:innen]
	\end{itemize}
	
	\subsection{Ausschusssitzung (\timeAuss)}
	\begin{enumerate}
		\item Präsentation: Ausschuss und Ablauf vorstellen \& erklären [5 min]
        \item Vorstellen der ÄA [10 min] \newline
        für jeden ÄA:
        \begin{enumerate}
            \item kurze Präsentation ÄA durch Ausschusssprecher:in
			\item kurze Bedenkzeit, dann Reaktion der anderen Fraktionen (startend links von vorstellender Fraktion $\to$ Uhrzeigersinn)
            \item mglw. Erwiderung einbringende Fraktion/weitere Diskussion \newline (Redezeitbegrenzung [bspw. 1 Min.] sinnvoll)
        \end{enumerate}
		\item Informelle Verhandlungen [40 min]
		\begin{itemize}
			\item Mehrheitsfindung/Koalitionen schmieden
			\item Kompromissfindung $\to$ ÄA können noch abgeändert/zusammengelegt werden, um mehrheitsfähig zu sein
		\end{itemize}
		\item Abstimmung über ÄA [20 min]\newline
		Für jeden ÄA:
		\begin{itemize}
			\item Abstimmung
			\item Ergebnis eintragen in Antragsgrün \newline
            $\to$ Kommentar verfassen mit \textit{Angenommen} oder \textit{Abgelehnt}
            (Angenommen, wenn mehr JA als NEIN Stimmen \newline [Enthaltungen zählen nicht])
		\end{itemize}
	\end{enumerate}
	
	\subsection{Mittagspause (\timeMittag)}
	\begin{itemize}
		\item Schüler:innen zu Plenarsaal bringen [Fraktionsteamer:innen]
		\item ÄA löschen/annehmen $\to$ Version 2 als \textit{Beschluss} anlegen und auf \textit{Eingereicht} stellen [\EP\ \& \KOM]
	\end{itemize}
	
	\subsection{2. Fraktionssitzung (\timeFrakTwo)}
	\begin{enumerate}
		\item Bericht der Schüler:innen aus ihren Ausschüssen [10 min]
		\item Finalen ÄA überlegen \& in Antragsgrün einbringen [10 min] \newline
		(\textbf{!!! in Version 2 arbeiten !!!}) \newline
		Bspw.:
		\begin{itemize}
			\item Zurückgewiesener ÄA in Ausschuss (mglw. abgeschwächt)
			\item Rücknahme eines im Ausschuss beschlossenen ÄA
			\item Überschüssige Idee aus 1. Sitzung \newline (eine Ausschussgruppe hatte mehr als eine ÄA-Idee)
		\end{itemize}
		
		\item Wahl Fraktionsvorsitz [5 min]
		\begin{itemize}
			\item Vorsitz: Abschlussrede (ca. 1 Min.)
			\item Vize: Vorstellung (ca. 30 Sek.) und Reaktion (ca. 10 Sek.) \newline bei finalen ÄA 
			\item Namen in Worddokument (Fraktionen/Fraktionsvorsitzende) \newline eintragen
		\end{itemize}
		\item Einteilung in Verhandlungsgruppen [5 min]
		\begin{itemize}
			\item Vier gleichgroße Gruppen (mehr oder weniger nach Wünschen der SuS)
			\item Auf andere Fraktionen aufteilen: Letzte Kompromisse mit dieser aushandeln
		\end{itemize}
	\end{enumerate}
	
	\subsection{Finale Verhandlungsphase (\timeFinVerh)}
	\begin{itemize}
		\item \stadtvertreter\ in Empfang nehmen [\KOM]
%		\item \politiker\ in Empfang nehmen [\KOM]

		\item \politiker\ in Empfang nehmen [Hannah]
		
		\item gegebenenfalls Fraktionsvorstand bei Rede unterstützen [Fraktionsteamer:innen]
		\item Währenddessen finale Verhandlungen zwischen Fraktionen über Kompromisse [Fraktionsteamer:innen]
		\begin{itemize}
			\item SuS in Plenarsaal bringen
			\item SuS unterstützen bei Kompromissformulierungen
			\item ÄA-Anpassungen bis 14.10 Uhr möglich (bei \EP\ \& \KOM\ melden)
		\end{itemize}
		\item ÄA in Version 1 übertragen \& Reihenfolge festlegen (weitreichendste zuerst) [\EP \& \KOM]
	\end{itemize}
	
	\subsection{Plenarsitzung (\timePlenar)}
	\begin{enumerate}
		\item Grußwort \newline [\stadtvertreter]
		\item Grußwort [\politiker]
		\item Plenardebatte [\EP]
		\begin{enumerate}
			\item Für jeden ÄA
			\begin{itemize}
				\item Vorstellung ÄA durch Fraktionsvize
				\item Reaktion durch übrige Fraktionsvize (Uhrzeigersinn)
				\item Abstimmung
				\item Zählen und Eintragen in Dashboard \newline (Mehr JA als NEIN Stimmen)
			\end{itemize}
			\item Alle angenommenen ÄA \& fertigen Gesetzentwurf vorstellen
			\item Abschlussrede durch Fraktionsvorsitz
			\item Abstimmung Gesetzesentwurf
			\item Zählen und Eintragen in Dashboard (Mehr JA als NEIN Stimmen)
		\end{enumerate}
		\item Debriefing via Slido
		\item Feedback über Jotform
		\item Vorstellung KV
		\item Danksagung an
		\begin{itemize}
            \item \stadtvertreter\ (\stadtvertreterOffice)
			\item \politiker\ (\politikerOffice)
			\item \stripPrefix{\location}
			\item \stripPrefix{\sponsor} \& \stripPrefix{\localSupport}
		\end{itemize}
	\end{enumerate}
	
	\subsection{Nachbereitung}
	\begin{itemize}
		\item Aufräumen
		\item Gruppenfoto
		\item Feedbackrunde
		\item Bericht (Website) \& Insta-Post
		\item Dankesmail an \stadtvertreter
		\item Dankesmail an \politiker
		\item Dankesmail an Lehrkräfte
	\end{itemize}
    \newpage
    \section{Packliste}
    \begin{itemize}
        \item Armbänder Fotoaufnahmen
        \item EU- \& JEF-Flaggen
        \item Klebeband, Schere \& Stifte
        \item Lanyards \& Hüllen
        \item Mappen
        \begin{itemize}
			\item (Flyer)
			\item Fraktionspapiere
			\item Gesetzesentwürfe
			\item Länderpapiere
        \end{itemize}
        \item Merch
        \item Namensschilder
        \item Raumschilder
        \item Teilnahmezertifikate
    \end{itemize}
    
    
    
    
    
    
    
    
    
\begin{comment}
    \newpage
    
    \section{Zeitstrahl}
    \begin{itemize}
        \item Im Jahr davor: Finanzierung
        \item Im Jahr davor: Ortfestlegung (Räumlichkeiten müssen stehen)
        \item Je nach Finanzierung: Terminfestlegung
        \item Direkt danach (mindestens halbes Jahr zuvor)
        \begin{itemize}
            \item Gastredner:innen anfragen
            \item Teamer:innen anfragen
            \item Schulklassen anfragen
            \item Hotel anfragen (mglw. Outsourcen)
        \end{itemize}
        \item 2-3 Monate davor: fixe Zusage der Teamer:innen einholen
        \item Direkt danach: Hotel buchen
        \item 1 Monat davor:
        \begin{itemize}
            \item Reminder an Lehrkräfte
            \item Catering abklären (mglw. Outsourcen)
            \item Aktualisierung der Unterlagen abgeschlossen
            \item mglw. Online-Schulung anbieten
        \end{itemize}
        \item 2 Wochen davor: 
        \begin{itemize}
            \item Restaurant reservieren (mglw. Outsourcen)
            \item KV Vorstellung abklären
        \end{itemize}
        \item In der Woche davor:
        \begin{itemize}
            \item Unterlagen vorbereiten
            \item Pressemitteilung rausschicken (mglw. Outsourcen)
        \end{itemize}
        \item Nach der SimEP: Nachbereitung
    \end{itemize}
\end{comment}
\end{document}